% Propietario conservado: /Users/ruben/Documents/New project/output/EXTREMA_MEDIA_RAZON_RECIPROCIDAD_ES_EN_20260918/ARTICULO/ES/source/sections/registro_k.tex
\section{Dodecaphasic extraction of the arithmetic--geometric coupling constant}
\label{x_reciprocidad:sec:k-registro}

\begin{definition}[Holographic arithmetic--geometric coupling constant]
\label{x_reciprocidad:lib01:definicion}
This is the name given to the ordered finite register produced by the oriented
extraction of the APP--TRIT--TPK state, with its phase origin, normalization and
declared readers, whose positional representation participates in the
arithmetic closure and whose incidence and geometric realizations act
on the same register. The symbol \(K\) designates its canonical integral
representation; \(\kappa_{\rm ag}\) designates its periodic rational
reading in the fixed chart.
\end{definition}

The generation of the regional coordinates does not exhaust the information of the
APP--TRIT--TPK state. An Archimedean reading determines a value through
compatible cylinders; another reading preserves the opposition of the sheets, the
orientation of the route and the register of incidences traversed. The
dodecaphasic register belongs to this second reading before participating in
the closure of alpha. Thus it is not defined by subtracting a physical constant from
a combination of previously known values. It is reconstructed from a
signed coordinate of the same arithmetic state with trajectory memory that produces the regional
coordinates.

The distinction is necessary even when two objects have twelve
coordinates. A word of twelve catalog blocks, a list of twelve
oriented events, a signed integer vector and a dodecaphasic register in base one thousand are
distinct objects. Comparing them requires the maps that
relate them. Equality of length does not provide those maps, and the
loss of signs is not merely a simplification of notation.

This chapter maintains the causal cut fixed in the preceding chapters:
the nine APP positions of the horizontal and vertical axes produce the
arithmetic sheets; TRIT types regime and orientation; TPK selects,
transports, updates and prolongs, preserving residue, quotient, carry,
sheet, route, boundary and memory. From the joint discrete structure of the
continuum one obtains the coordinates used here. The chain
\(w_6\to w_{12}\to w_{18}\to w_{24}\to w_{30}\to R_{36}\to G_9\)
is not replaced by a finite calendar. Nor is the order
\(U_t\to\Gamma_9\to K_{\mathrm{ph}}\) reversed: the final map is a
reduced memory reading, not the generator of the signed register.


% BEGIN OWNER /Users/ruben/Documents/New project/output/EXTREMA_MEDIA_RAZON_RECIPROCIDAD_ES_EN_20260918/ARTICULO/ES/source/sections/revision_k.tex
% Adición propuesta al inicio de registro_k.tex, después de su introducción.
% Conserva todas las definiciones, fórmulas y demostraciones actuales.
% Procedencia: registro_k.tex, §§k-hadamard/k-transversal/k-kappas;
% k_reversibilidad.tex; incidencia_articulacion.tex, inc-ampl-cuaterna.
\paragraph{Reversible coordinatization and incidence determination.}
\label{x_reciprocidad:ampl-alfa-k-intro-reversible}
The result organizing this chapter is the equivalence between three
coordinatizations of a register previously produced by
APP--TRIT--TPK. The Hadamard transformation relates the signed
register to the integer register; the difference extractor relates
the latter to its transverse variations and its uniform component:
\[
 U\underset{\mathcal H_{12}}{\overset{\mathcal H_{12}/4}{\longleftrightarrow}}K
 \overset{(D_3,D_4,Q)}{\longleftrightarrow}
 (D_3K,D_4K,Q(K)).
\]
The first equivalence is established on the integral sublattice
$\mathcal L_H$; the second, on the image of the extractor. With the
base $1000$, the $12$ positions and the phase origin fixed, the evaluation
$K\mapsto\kappa_{\mathrm{per}}$ adds a fourth reversible
coordinatization: multiplying by $1000^{12}-1$ and performing Euclidean
divisions recovers the complete register. The proofs of
\S\ref{x_reciprocidad:subsec:k-hadamard}, \S\ref{x_reciprocidad:subsec:k-transversal} and
\S\ref{x_reciprocidad:subsec:k-kappas} establish these equivalences in their respective
domains.

Two compositions then act on that register. Positional
composition combines its components with the closure,
propagation and autoscale coordinates, and its normalization determines $\alpha_A$. Incidence
composition applies the cubic threshold and obtains
\[
 K\longmapsto B_K=\{j:K_j\geq729\}=\{4,6,9,10\}.
\]
The complete register allows both compositions to be reconstructed. The support
$B_K$ specifically preserves membership in the class selected
by the threshold. The identity $B_K=H_e\cap P_\pi$ of
\eqref{x_reciprocidad:inc-ampl-cuaterna} links this integer extraction to the incidence
of the regional codes. The negative support of its balance with $K$
completes the flag $B_K\subset H^-_\alpha$ developed in the
exceptional prolongation. Thus the function of $K$ is defined through
an invertible transformation and through precise subsequent
maps; the additional depth and transport information remains
in the state from which the register originates.

% END OWNER


\subsection{Cyclic extension and memory grading}
\label{x_reciprocidad:subsec:k-calendario}

The observable calendar contains one hundred and eight positions, organized into
twelve nonadic windows. With positive indices, let
\[
 E_m=\{9(m-1)+1,\ldots,9m\},\qquad 1\leq m\leq12.
\]
The ordered support \(E_1\sqcup\cdots\sqcup E_{12}\) preserves position and
window membership. The group \(C_{108}=\mathbb Z/108\mathbb Z\)
instead preserves the cyclic translation law. For the representative
\(0\leq\widetilde\theta<108\), the coordinates
\[
 \beta(\theta)=\left[\left\lfloor\frac{\widetilde\theta}{9}\right\rfloor\right]_{12},
 \qquad r(\theta)=1+(\widetilde\theta\bmod9)
\]
separate the dodecaphasic sector and the interior nonadic position. Neither
by itself preserves the total number of turns.

\begin{proposition}[The calendar carry]
\label{x_reciprocidad:prop:k-extension}
The maps
\[
 \iota:C_{12}\longrightarrow C_{108},\quad[b]\longmapsto[9b],
 \qquad p:C_{108}\longrightarrow C_9,\quad[\theta]\longmapsto[\theta]_9
\]
form a nonsplit exact sequence
\[
 0\longrightarrow C_{12}\xrightarrow{\iota}C_{108}
 \xrightarrow{p}C_9\longrightarrow0.
\]
For the set-theoretic section choosing the representatives
\(0,\ldots,8\), the additivity defect is the carry
\[
 c(r,r')=\left[\left\lfloor\frac{\widetilde r+\widetilde r'}9\right\rfloor\right]_{12}.
\]
\end{proposition}

\begin{proof}
The kernel of \(p\) consists of the multiples of nine and coincides with
the image of \(\iota\). Reduction is surjective, and multiplication
by nine is injective on the stated domain. If the sequence
split, the middle group would be isomorphic to \(C_{12}\times C_9\), whose
exponent is \(\operatorname{mcm}(12,9)=36\); the exponent of
\(C_{108}\) is 108. This contradiction proves nonsplitting. Finally,
Euclidean division of \(\widetilde r+\widetilde r'\) by nine gives the
stated carry and verifies the identity for the section.
\end{proof}

The result explains a difference that will reappear in alpha: returning to the
phase is not equivalent to returning to the state. If \(q_{\mathrm{mem}}\) counts the
nonadic turns, the calendar preserves only
\(\beta=[q_{\mathrm{mem}}]_{12}\). After one hundred and eight steps,
\(\beta\) returns, but \(q_{\mathrm{mem}}\) increases by twelve. The
register also preserves the cylinder, the incidence register and the vertical
memory. The closure of a finite chart does not turn that history into a
memoryless cycle.

\subsection{The oriented register and the integral chart \texorpdfstring{\(1+5+6\)}{1+5+6}}
\label{x_reciprocidad:subsec:k-registro-integral}

For a history produced by TPK, let \(\tau_m\) be the subroute over
\(E_m\), \(\sigma_m\) its orientation, \(\kappa_m\) the boundary
carry, and \(\Lambda_m\) the local incidence entry. The register is
\[
 \mathcal R_{12}(x)=
 \bigl((E_m,\tau_m,\sigma_m,\kappa_m,\Lambda_m)\bigr)_{m=1}^{12}.
\]
Its domain is the space of arithmetic histories with trajectory memory, not the set of
phase labels. Projection onto the calendar forgets precisely some
of the data that the signed reader needs.

An initial chart of the register can be written through oriented
events. Oriented extraction fixes the set of event codes
\(\{2,4,6,8\}\) and the sectoral orientation
\[
 \Sigma_{12}=(+,+,+,-,-,-,+,+,+,-,-,-).
\]
If \(d_t\) is the direction code at step \(t\), one obtains
\[
 e_i=\Sigma_{12}(i+1)
 \#\{t\in E_{i+1}:d_t\in\{2,4,6,8\}\},
 \qquad0\leq i<12.
\]
The formula keeps the nonnegative census of events and the
orientation separate. It does not identify a cardinal direction with a TRIT sign. The
vector \(e\), whose components are signed local censuses, is not yet
the vector \(U\) that will appear below either.

The opposite positions determine, for \(0\leq i<6\),
\[
 p_i=e_i+e_{i+6},\qquad a_i=e_i-e_{i+6}.
\]
With
\[
 s_C=\sum_{i=0}^{5}p_i,\qquad M_i=p_i-p_5\quad(0\leq i<5),
\]
one defines the integral chart
\[
 \mathcal L_{12}(e)
 =(s_C,M_0,\ldots,M_4,a_0,\ldots,a_5).
\]
The decomposition \(1+5+6\) is not an arbitrary distribution of coordinates:
one coordinate preserves the charge, five compare the pairs relative to a
reference pair and six preserve the opposition of the half-crowns.

\begin{lemma}[Inversion of the integral chart]
\label{x_reciprocidad:lem:k-carta-integral}
An integer tuple \((s_C,M_0,\ldots,M_4,a_0,\ldots,a_5)\) belongs to the
image of \(\mathcal L_{12}\) if and only if
\[
 s_C-\sum_{i=0}^{4}M_i\equiv0\pmod6,
 \qquad p_i\equiv a_i\pmod2\quad(0\leq i<6),
\]
where
\[
 p_5=\frac{s_C-\sum_{i=0}^{4}M_i}{6},\qquad p_i=p_5+M_i\quad(i<5).
\]
In that case its preimage is unique and is given by
\[
 e_i=\frac{p_i+a_i}{2},\qquad e_{i+6}=\frac{p_i-a_i}{2}.
\]
\end{lemma}

\begin{proof}
The sum of the five margins is \(s_C-6p_5\), from which division
by six is obtained. Once the \(p_i\) are reconstructed, each pair of equations
\(p_i=e_i+e_{i+6}\), \(a_i=e_i-e_{i+6}\) has the indicated solution.
This solution is integral exactly under the parity congruence. All
the operations are inverted, so no residual freedom remains.
\end{proof}

This reversibility shows what information is preserved; it does not authorize
removing the rest of the article. The complete register has a larger domain
than this census chart. Extracting the channels that feed
\(U\) uses that complete register and does not follow from the list
\((e_i)\) through an identification of names.


% BEGIN OWNER /Users/ruben/Documents/New project/output/EXTREMA_MEDIA_RAZON_RECIPROCIDAD_ES_EN_20260918/ARTICULO/ES/source/sections/registro_incidencias.tex
% Suplemento de trabajo. No sustituye registro_k.tex ni cierra por decreto E108.
% Procedencia y alcance: INFORME_RECUPERACION.md, misma carpeta.
\subsection{Incidence evaluation preceding the dodecaphasic register}
\label{x_reciprocidad:subsec:registro-evaluacion-incidencial}

Recovering a register from its cyclic differences and
total charge is a different question from evaluating those coordinates
on the incidence register. The first problem is solved by the
transverse extractor of Proposition~\ref{x_reciprocidad:prop:k-transversal}.
For the second, a historical presentation of the
register provides three integer counts in each window. It is useful
to make their composition explicit and determine exactly the information
required for their evaluation.

Let \(\mathscr L\) be a previously generated ledger of visits and oriented
traces. Each visit preserves the route, the APP position, the sheet, the
instant, the multiplicity and the boundary incidence. The arithmetic
evaluation produces the integer and its decomposition
\[
 n_t=r_t+9q_t,\qquad 1\leq r_t\leq9.
\]
The integer \(n_t\) is calculated from the sum or product sheet; its
residue is not received as an independent value. The distinction between
crown incidence and interior incidence remains explicit for
visits of residue \(9\).

In window \(E_m\), the count presentation considers
\begin{align}
 A_m&=\sum_{v\in E_m}w(v)
             \mathbf1_{\{1,3,5,7\}}(r(v)),\\
 C_m&=\sum_{j\geq0}\bigl(\nu^-_{120,m}(j)-\nu^+_{120,m}(j)\bigr),\\
 V_m&=\sum_{v\in E_m}w(v)\mathbf1_{\{9\}}(r(v))\epsilon_\partial(v).
\end{align}
Here \(w(v)\) is the incidence multiplicity, and
\(\epsilon_\partial(v)=1\) on the corona, \(-1\) in the interior.
The traces counted by \(\nu^\pm\) must be enumerated by the
corresponding rule in the article. The integer sum requires finite support or an
alternative construction determining its meaning. The number of
chronological incidences and the reduced length \(120(1+3j)\) have
different types; relating them requires the reader's unit of length.

One defines \([x]_{10}\in\{0,\ldots,9\}\) as the Euclidean
representative and preserves the quotients
\[
 x=[x]_{10}+10\lfloor x/10\rfloor.
\]
The incidence block is
\begin{equation}
 z_m=100[A_m]_{10}+10[C_m]_{10}+[V_m]_{10}.
 \label{x_reciprocidad:eq:registro-bloque-incidencial}
\end{equation}
Composition yields the map
\begin{equation}
 \mathcal E_{\rm cnt}(\mathscr L)=
 \left((z_m-z_{m+3})_{m=1}^{12},
       (z_m-z_{m+4})_{m=1}^{12},\sum_{m=1}^{12}z_m\right).
\label{x_reciprocidad:eq:registro-composicion-incidencial}
\end{equation}
The indices are read cyclically modulo \(12\). In this subsection,
\((D_jz)_m=z_m-z_{m+j}\), for \(j=3,4\), and
\(Q(z)=\sum_{m=1}^{12}z_m\); therefore, the preceding composition is
\((D_3z,D_4z,Q(z))\). This convention is preserved in
\S\ref{x_reciprocidad:subsec:k-transversal}, where complete integral
reconstruction is proved.
This formula produces its coordinates from the counts. It does not use
\(K\), \(U\), \(\alpha\) or the expected transverse coordinates
as arguments. Its identification with
\(\mathcal E_{108}^{90,120}\) requires composing it with the family and
the effective incidences of the terminal state; it does not follow solely
from the invertibility of the transverse system.

\begin{proposition}[Sufficient and necessary arithmetic information]
\label{x_reciprocidad:prop:registro-36-residuos}
Let \(N,N'\in\mathbb Z^{12\times3}\), ordered by the counts
\((A,C,V)\). For the map defined by
\eqref{x_reciprocidad:eq:registro-bloque-incidencial}--\eqref{x_reciprocidad:eq:registro-composicion-incidencial},
\[
 \mathcal E_{\rm cnt}(N)=\mathcal E_{\rm cnt}(N')
 \quad\Longleftrightarrow\quad
 N-N'\in10\mathbb Z^{36}.
\]
In particular, the \(36\) sectoral residues exactly determine
the output of \(25\) coordinates. Preserving the quotients and
the provenance constitutes additional information of the register.
\end{proposition}

\begin{proof}
Equality modulo \(10\) produces the same blocks and therefore
the same transverse coordinates. Conversely, if the transverse
coordinates coincide, \(D_3(z-z')=D_4(z-z')=0\). The steps
\(3\) and \(4\) generate the cycle of length \(12\), so
\(z-z'\) is constant. Equality of total charges cancels that
constant. Finally, the decimal representation of an integer between
\(0\) and \(999\) has three unique digits. This recovers the
three residues of each window. The assertion is an algebraic identity
for every \(N,N'\), not an inference obtained from finite tests.
\end{proof}

\subsection{Conjugation of counts and boundary terms}
\label{x_reciprocidad:subsec:registro-conjugacion-incidencial}

Let \(\rho(m)=13-m\). Consider a conjugation of incidences
that reverses the order of the windows, preserves the arithmetic class and the
boundary class of the visits and interchanges the channels of the traces.
Then
\[
 (A_m^\dagger,C_m^\dagger,V_m^\dagger)
 =(A_{\rho(m)},-C_{\rho(m)},V_{\rho(m)}).
\]
If \(c_m=[C_m]_{10}\), decimal reduction produces the identity
\begin{equation}
 z_m^\dagger=z_{\rho(m)}
       +100\mathbf1_{\{c_{\rho(m)}\ne0\}}-20c_{\rho(m)}.
 \label{x_reciprocidad:eq:registro-conjugacion-decimal}
\end{equation}
Indeed, \([-C]_{10}=0\) when \([C]_{10}=0\), and
\([-C]_{10}=10-[C]_{10}\) otherwise. The formula preserves
exactly that difference. Negating the channel is not equivalent to negating
the entire decimal block.

Applying this identity to a concrete involution of TPK requires
verifying its action on the incidences. For an occupancy reader
evaluated before the advance, reversal of a route
\(\gamma=(x_0,\ldots,x_n)\) satisfies
\[
 \sum_{k=0}^{n-1}f(x_{n-k})
 =\sum_{k=0}^{n-1}f(x_k)+f(x_n)-f(x_0).
\]
The endpoint terms are incorporated before reducing modulo
\(10\). In closed routes they vanish; in open windows they may
alter their residues. This correction, already present in the bidirectional
development of transport, specifies what the junction
between the route and the counts must preserve.

\begin{proposition}[Necessity of the nonperiodic information of the register]
Suppose that the observable of a fixed family of routes satisfies
\(o(t+54)=o(t)\), that multiplicity is preserved under that return
and that the same local reader acts in the windows of \(9\) instants.
Then its counts and blocks satisfy
\[
 (A,C,V)_{m+6}=(A,C,V)_m,\qquad z_{m+6}=z_m.
\]
\end{proposition}

\begin{proof}
The translation \(t\mapsto t+54\) takes \(E_m\) bijectively
to \(E_{m+6}\), preserving each counted incidence and its multiplicity.
Equality is transmitted to the sum and to decimal reduction.
\end{proof}

Therefore, a presentation with \(12\) distinct blocks requires the
reader to preserve some information that does not factor through that observable
of period \(54\): incidence selection, effective orientation,
endpoints, multiplicity or memory. The proposition delimits an inadmissible
reduction of the state; it does not state that the full dynamics has
period \(54\) or determine by itself what its terminal reader is.

% END OWNER


% BEGIN OWNER /Users/ruben/Documents/New project/output/EXTREMA_MEDIA_RAZON_RECIPROCIDAD_ES_EN_20260918/ARTICULO/ES/source/sections/censo_terminal_completo.tex
\subsection{Complete census of cylinders and selection of the terminal boundary}
\label{x_reciprocidad:censo:terminal-completo}

The terminal reading jointly uses two pieces of information: compatibility
of the regional cylinders and an incidence profile with labeled channels and
positions. This subsection constructs the complete catalogs,
enumerates their triples and proves the effect of each condition of the reader.
The starting point is the regional outputs already produced in
section~\ref{x_reciprocidad:sec:generacion}, not values introduced to define
APP--TRIT--TPK transport. Reduction to a matrix belongs to a subsequent
reading of those outputs and preserves the order of the three channels.

The result is a census theorem with an explicit starting structure.
The profile of Gram, weights, distances and margins is declared as a datum of the
terminal reader. The exhaustiveness of the enumeration proves what
that profile selects on the fixed catalogs; it does not by itself prove that the
profile is a universal law or that it can be reconstructed from the unique triple
that ultimately satisfies it. This separation allows the complete finite
proof to be incorporated without reversing its genealogy.

\subsubsection{From regional outputs to ternary catalogues}

Let \(x_\pi,x_e,x_\varphi\in[0,1)\) be the fractional parts of the
three regional coordinates already constructed. Their published cylinders
of one thousand decimal positions are used,
\[
 J_\chi^{(1000)}=
 \left[\frac{d_\chi}{10^{1000}},
       \frac{d_\chi+1}{10^{1000}}\right),
 \qquad \chi\in\{\pi,e,\varphi\}.
\]
The integers \(d_\chi\) are finite output readings. The length one thousand
is sufficient for the integers that follow, as is checked through the
two bounds of each interval; it is not a depth assumed infinite.
We define
\begin{equation}
 p=30,\quad r=1050,\quad L=p+r=1080,\quad k=171,
 \qquad D=3^L,\quad Q=1000^k,
 \label{x_reciprocidad:censo:parametros}
\end{equation}
and
\[
 P_\chi=\lfloor3^{30}x_\chi\rfloor,\qquad
 T_\chi=\lfloor Qx_\chi\rfloor.
\]
Here \(k\) counts decimal triads and does not designate the dodecaphasic register
\(K\). One has \(1000^{171}\leq3^{1080}<1000^{172}\).

\begin{lemma}[Stable extraction of an integer]
\label{x_reciprocidad:censo:estabilidad}
For \(d,m,h\in\mathbb Z\), with \(m,h>0\), the value of
\(\lfloor mx\rfloor\) is constant on \([d/h,(d+1)/h)\) if and only if
\begin{equation}
 \left\lfloor\frac{md}{h}\right\rfloor
 =\left\lfloor\frac{m(d+1)-1}{h}\right\rfloor.
 \label{x_reciprocidad:censo:prueba-estabilidad}
\end{equation}
When this holds, either side gives that value.
\end{lemma}
\begin{proof}
The minimum of \(mx\) is attained at the left endpoint. Its supremum
is \(m(d+1)/h\) and is not attained. The largest integer that
\(\lfloor mx\rfloor\) can take is
\(\lceil m(d+1)/h\rceil-1\), equal to the right-hand side because
the numerator and denominator are integers. Equality of the endpoints
of the value interval proves the assertion.
\end{proof}

The stability proof holds in all three channels for
\(m=3^{30},Q,3^{1080}\). The last choice will be used only
for a subsequent comparison, not to select the profile. The
first thirty recovered positions are precisely the five
stages of six trits of regional transport:
\begin{center}\small
\begin{tabular}{@{}c l r@{}}\toprule
Channel&Thirty-trit prefix&\(P_\chi\)\\\midrule
\(\pi\)&\texttt{010211012222010211002111110221}&29152671743887\\
\(e\)&\texttt{201101121221102011012222102011}&147887858824447\\
\(\varphi\)&\texttt{121200112202121020010210010200}&127247717616687\\\bottomrule
\end{tabular}
\end{center}
The coincidence preserves the prolongation already constructed; the decimal reading
is not used to retroactively replace its transport matrices.

For a fixed channel, a tail of \(r\) trits is represented by
\(0\leq R<3^r\). Its cylinder and the cylinder of the first \(k\)
triads are
\begin{equation}
 I_{\chi,R}=\left[\frac{P_\chi3^r+R}{D},
                       \frac{P_\chi3^r+R+1}{D}\right),\qquad
 J_{\chi,T}=\left[\frac{T_\chi}{Q},\frac{T_\chi+1}{Q}\right).
 \label{x_reciprocidad:censo:dos-cilindros}
\end{equation}
The catalog includes all tails for which these two cylinders
have nonempty intersection.

\begin{proposition}[Catalogue by exact intersection]
\label{x_reciprocidad:censo:catalogos}
The compatible tails are the integers in the interval \([a_\chi,b_\chi]\),
where
\begin{align}
 a_\chi&=\max\left(0,
       \left\lfloor\frac{T_\chi D}{Q}\right\rfloor-P_\chi3^r\right),
       \label{x_reciprocidad:censo:extremo-a}\\
 b_\chi&=\min\left(3^r-1,
       \left\lfloor\frac{(T_\chi+1)D-1}{Q}\right\rfloor-P_\chi3^r\right).
       \label{x_reciprocidad:censo:extremo-b}
\end{align}
For the fixed regional cylinders, the truncations do not act and one obtains
\begin{equation}
\begin{array}{c|r|r|r}
 \chi&a_\chi\bmod729&b_\chi\bmod729&b_\chi-a_\chi+1\\\hline
 \pi&67&262&196\\
 e&584&51&197\\
 \varphi&117&312&196
\end{array}
\label{x_reciprocidad:censo:tabla-catalogos}
\end{equation}
\end{proposition}
\begin{proof}
Write \(N=P_\chi3^r+R\). Two half-open intervals of
\eqref{x_reciprocidad:censo:dos-cilindros} intersect exactly when
\[
 NQ<(T_\chi+1)D,\qquad (N+1)Q>T_\chi D.
\]
The second inequality is equivalent to
\(N\geq\lfloor T_\chi D/Q\rfloor\); the first, to
\(N\leq\lfloor((T_\chi+1)D-1)/Q\rfloor\). Subtracting
\(P_\chi3^r\) and preserving \(0\leq R<3^r\) proves the formulas.
Evaluation by integer divisions of the \(T_\chi\) extracted by
\eqref{x_reciprocidad:censo:prueba-estabilidad} gives the table. Since \(r\geq6\),
\(P_\chi3^r\equiv0\pmod{729}\), so the residue of the first
endpoint is also obtained directly from
\(\lfloor T_\chi D/Q\rfloor\bmod729\).
\end{proof}

The last six trits of a tail form the word
\(\nu_6(R\bmod729)\), where \(\nu_6\) is the ternary notation
of length six, including its leading zeros. By the table, each
catalog has fewer than \(729\) consecutive elements, so
this reduction is injective on it. Thus the three boundary catalogs
are completely described, without choosing the final triples:
\begin{equation}
 \begin{split}
 \mathcal B_\pi&=\{\nu_6(67+i):0\leq i<196\},\\
 \mathcal B_e&=\{\nu_6((584+j)\bmod729):0\leq j<197\},\\
 \mathcal B_\varphi&=\{\nu_6(117+\ell):0\leq\ell<196\}.
 \end{split}
 \label{x_reciprocidad:censo:fronteras-explicitas}
\end{equation}
Requiring the left endpoint of \(I_{\chi,R}\) to belong to
\(J_{\chi,T}\) would be another condition: it would lose the first cylinder
that intersects the boundary and give \(195,196,195\). The census uses the
intersection of the complete cylinders, in accordance with its definition.

\subsubsection{Declared profile and compatibility indicators}

For \(u,v\in\mathbb F_3^6\), let
\(g(u,v)=\sum u_iv_i\in\mathbb F_3\),
\(n(u)=g(u,u)\), \(w(u)=\#\{i:u_i\ne0\}\) and
\(h(u,v)=\#\{i:u_i\ne v_i\}\). The last two functions take
integer values; the first two, values in \(\mathbb F_3\).
The local profile of the reader, in the order \((\pi,e,\varphi)\), is
\begin{equation}
 \begin{split}
 (g_{\pi e},g_{\pi\varphi},g_{e\varphi})&=(2,2,0),\qquad
 (n_\pi,n_e,n_\varphi)=(0,0,0),\\
 (w_\pi,w_e,w_\varphi)&=(3,3,3),\qquad
 (h_{\pi e},h_{\pi\varphi},h_{e\varphi})=(2,5,3).
 \end{split}
 \label{x_reciprocidad:censo:perfil-local}
\end{equation}
The words are ordered as in \eqref{x_reciprocidad:censo:fronteras-explicitas}.
We shall use \([E]\in\{0,1\}\) for the indicator of a condition \(E\).
For each pair of channels \(\chi,\psi\), the rectangular matrix
\(A_{\chi\psi}\) has entry \([g(u,v)=g_{\chi\psi}]\).
The matrix \(A^{h}_{\chi\psi}\) multiplies that entry by
\([h(u,v)=h_{\chi\psi}]\). Let \(D_\chi^n\) be the diagonal matrix
of \([n(u)=0]\), and \(D_\chi^{nw}\) that of
\([n(u)=0][w(u)=3]\).

\begin{theorem}[Exhaustive enumeration by finite sums]
\label{x_reciprocidad:censo:teorema-conteo}
The product of the three catalogues contains \(7\,567\,952\) triples.
The successive conditions in \eqref{x_reciprocidad:censo:perfil-local} leave
\begin{equation}
 7\,567\,952\longrightarrow275\,794\longrightarrow6235
 \longrightarrow3127\longrightarrow331.
 \label{x_reciprocidad:censo:cadena-local}
\end{equation}
These four counts are, respectively,
\begin{align}
 N_g&=\operatorname{tr}(A_{\pi e}A_{e\varphi}A_{\varphi\pi}),
 \nonumber\\
 N_{gn}&=\operatorname{tr}(D_\pi^n A_{\pi e}
             D_e^n A_{e\varphi}D_\varphi^n A_{\varphi\pi}),
 \nonumber\\
 N_{gnw}&=\operatorname{tr}(D_\pi^{nw} A_{\pi e}
             D_e^{nw} A_{e\varphi}D_\varphi^{nw} A_{\varphi\pi}),
 \nonumber\\
 N_{gnwh}&=\operatorname{tr}(D_\pi^{nw} A^h_{\pi e}
             D_e^{nw} A^h_{e\varphi}D_\varphi^{nw} A^h_{\varphi\pi}).
 \label{x_reciprocidad:censo:trazas}
\end{align}
All products and traces in this formula are calculated in the integers,
not in the field of three elements.
\end{theorem}
\begin{proof}
The size of the product is \(196\cdot197\cdot196\). When a
trace of \eqref{x_reciprocidad:censo:trazas} is expanded, its indices run through exactly one
triple of words. The summand equals one if it satisfies all the represented
conditions and zero otherwise. Each triple is counted once.
The diagonal matrices add the individual conditions; the rectangular
factors add the pairwise conditions. This proves that the
traces are precisely the stated cardinalities.

To make their evaluation explicit, once the indices \(i,j\) of the first
two channels are fixed, one forms the sets of indices \(\ell\) of the third
that satisfy each pairwise condition. Their indicators are intersected,
the relevant individual indicators are added and their cardinality is summed.
The final sum runs over \(0\leq i<196\), \(0\leq j<197\).
Substituting the words of \eqref{x_reciprocidad:censo:fronteras-explicitas} and
the profile \eqref{x_reciprocidad:censo:perfil-local} gives, in that order,
\(275794,6235,3127,331\). The supplement records the four contributions
of each of the \(196\) indices \(i\) and the \(331\) complete triples;
its program evaluates these sums with exact binary indicators and checks
each survivor scalarly. Neither a sample nor a tolerance is involved.
\end{proof}

\subsubsection{Labelled incidence and orientation charge}

For a surviving triple, let \(B\in\mathbb F_3^{3\times6}\) be its
matrix of ordered rows. The second profile requires
\begin{equation}
 \operatorname{rank}_{\mathbb F_3}B=3,\qquad
 w_{\mathrm{col}}(B)=(1,1,2,2,3,0),\qquad
 q_{\partial}(B)=(1,2,2,1,2,0),
 \label{x_reciprocidad:censo:perfil-columnas}
\end{equation}
where \(w_{\mathrm{col}}\) counts the nonzero elements of each
column and \(q_{\partial}\) sums its entries modulo three. It is a profile
of labeled columns; permuting them without transporting the profile defines another
reader. Denote by \(\mathcal F\) the set of the \(331\) triples
of the preceding theorem.

\begin{proposition}[Reduction of the census to the terminal pair]
\label{x_reciprocidad:censo:dos-matrices}
The sums of indicators over \(\mathcal F\), successively requiring
rank, column weights and column sums, are
\begin{equation}
 \sum_{B\in\mathcal F}[\operatorname{rank}B=3]=331,
 \qquad N_{\mathrm{rango,pesos}}=18,
 \qquad N_{\mathrm{rango,pesos,sumas}}=2.
 \label{x_reciprocidad:censo:conteo-columnas}
\end{equation}
With catalogue indices starting at one, the two final triples are
\((120,197,157)\) and \((129,197,148)\), and their matrices are
\[
 B_0=\begin{pmatrix}0&2&0&2&2&0\\0&0&1&2&2&0\\1&0&1&0&1&0\end{pmatrix},
 \qquad
 B_1=\begin{pmatrix}0&2&1&0&2&0\\0&0&1&2&2&0\\1&0&0&2&1&0\end{pmatrix}.
\]
\end{proposition}
\begin{proof}
For each triple in the set already enumerated, the rank is calculated by elimination
in \(\mathbb F_3\), the nonzero elements of the six
columns are counted and their entries are summed modulo three. The products of the indicators
of \eqref{x_reciprocidad:censo:perfil-columnas} give \eqref{x_reciprocidad:censo:conteo-columnas}.
Substituting the two final indices into
\eqref{x_reciprocidad:censo:fronteras-explicitas} gives the displayed matrices. The complete
list of survivors and the exact filters in the supplement allow
the elimination to be reproduced without presupposing those two matrices as the domain.
\end{proof}

The orientation \(\sigma=(1,1,-1)=(1,1,2)\in\mathbb F_3^3\)
selects the line
\(\langle\sigma\rangle=\{(0,0,0),(1,1,2),(2,2,1)\}\).
The row charges of \(B_0,B_1\) are \((0,2,0)\) and \((2,2,1)\);
only the second belongs to that line. This recovers the last step
of the selector in section~\ref{x_reciprocidad:subsec:k-terminal}, now preceded by
the construction and enumeration of its complete domain. The stable reading
of \(\lfloor3^{1080}x_\chi\rfloor\), performed after the filters,
gives the boundaries \texttt{021020}, \texttt{001220}, \texttt{100210},
with the same ranks \((129,197,148)\). This subsequent comparison does not define
the profile or participate in the indicator sums.

\subsubsection{Preserved information and scope of the result}

The procedure preserves the thirty-trit prefix, the compatible tail,
the regional order, the position of each trit and the orientation. Reduction
modulo \(729\) here preserves the identity of each candidate because the
tail interval has length less than \(729\); this injectivity
is not asserted for an arbitrary history. The census of \(331\) also proves
that the isolated local profile does not select a single matrix: the
labeled margins and the charge contribute additional effective conditions.

Thus a complete finite enumeration and its selector
on output cylinders have been recovered, with an explicit reading profile. The construction
of the signed register additionally uses its memory reader and its channels,
as specified below; the visible matrix is not identified with
that register by equality of the number of components. Nor does this count
prove the universality of the profile, generate the regional coordinates
again or decide an assertion about alpha or about the complete continuum.

% END OWNER


\subsection{Two readings of the terminal state}
\label{x_reciprocidad:subsec:k-terminal}

Denote by \(x_{\mathrm{term}}\) the state produced by the composition
of selection, transport and update up to the terminal cut. The
operator notation is
\[
 x_{\mathrm{term}}=
 (\mathcal U_{108}\circ\cdots\circ\mathcal U_1)(x_{\mathrm{seed}}),
 \qquad \mathcal U_t=\mathsf{Upd}_t\circ\mathsf{Tra}_t\circ\mathsf{Sel}_t.
\]
Its two relevant readings are
\[
 x_{\mathrm{term}}
 \xrightarrow{(\pi_{\mathrm{term}},R_{\mathrm{sgn}})}
 (B_{\mathrm{term}},U).
\]
The matrix \(B_{\mathrm{term}}\) retains the visible ternary reading; the
vector \(U\) preserves another coordinate of the same history. No arrow
\(B_{\mathrm{term}}\to U\) is inserted between them.

The terminal selector operates on the three regional catalogues
of sizes 196, 197, and 196, constructed in
Section~\ref{x_reciprocidad:censo:terminal-completo}. The local signatures reduce their
\(196\cdot197\cdot196=7\,567\,952\) triples to 331, and the column
marginal \(q_{\partial}=(1,2,2,1,2,0)\) leaves two matrices:
\[
 B_0=\begin{pmatrix}0&2&0&2&2&0\\0&0&1&2&2&0\\1&0&1&0&1&0\end{pmatrix},
 \qquad
 B_1=\begin{pmatrix}0&2&1&0&2&0\\0&0&1&2&2&0\\1&0&0&2&1&0\end{pmatrix}.
\]
The row orientation is \(\sigma=(1,1,-1)=(1,1,2)\) in
\(\mathbb F_3^3\). The row charges of the candidates are
\((0,2,0)\) and \((2,2,1)\). Only the second belongs to
\(\langle\sigma\rangle=\{(0,0,0),(1,1,2),(2,2,1)\}\); therefore
\(B_{\mathrm{term}}=B_1\). This check proves the last step of the
selector from the declared census. It does not turn the column margin
into a consequence of the row charge, or replace the preceding enumeration
with the inspection of those two matrices.

The signed reader has the typed composition
\begin{equation}
 R_{\mathrm{sgn}}
 =\Pi_H\circ\mathcal E_{108}^{90,120}\circ\mathcal R_{12}:
 \mathcal X_{\mathrm{TPK}}^{\mathrm{term,mem}}
 \longrightarrow\mathbb Z^{12}.
 \label{x_reciprocidad:eq:k-lector-firmado}
\end{equation}
The channel register determines
\[
 \mathcal E_{108}^{90,120}(\mathcal R_{12}(x_{\mathrm{term}}))
 =(b^{90},b^{120},Q_{\mathrm{TPK}}),
\]
with coordinates
\begin{align*}
 b^{90}={}&(-495,-116,-684,108,601,-90,\\
            &\hspace{19mm}-173,-88,313,560,-397,461),\\
 b^{120}={}&(-425,-281,-481,671,-255,30,\\
             &\hspace{19mm}475,-543,680,251,6,-128),\\
 Q_{\mathrm{TPK}}={}&6263.
\end{align*}
These quantities are used as coordinates of the preceding register, not
as parameters adjusted through alpha. The reconstruction proved
in the following subsections begins exactly at this typed output.
The integral characterization of the extractor, its inverse and the composition
of contributions are developed in section~\ref{x_reciprocidad:img09:seccion},
in particular in Theorem~\ref{x_reciprocidad:img09:imagen} and
Proposition~\ref{x_reciprocidad:img09:acumulacion}. The regional determination
of the descriptor and terminal selection are presented in
chapter~\ref{chap:despliegue-selector-k}. Each construction preserves its
domain: the selector acts on the generated repertoires and inversion
acts on the signed channels declared here. The documentary trail of the
extractor is recorded in the technical provenance of the package.

\subsection{Construction of the signed register through the oriented channels}
\label{x_reciprocidad:subsec:k-pi-h}

The harmonic reader \(\Pi_H\) can be written without using either \(K\) or
alpha. To avoid confusing its orbital sums with the alpha blocks,
we shall denote them by \(s_1,s_2,s_3\). Let
\[
 B_r=\sum_{j=0}^{3}b^{120}_{r+3j},\qquad
 d_{r,j}=b^{90}_{r+3j},\qquad1\leq r\leq3,
\]
and
\begin{equation}
 s_1=\frac{Q_{\mathrm{TPK}}+2B_1+B_2}{3},\qquad
 s_2=s_1-B_1,\qquad s_3=s_2-B_2.
 \label{x_reciprocidad:eq:k-sumas-armonicas}
\end{equation}
Division by three belongs to the reader of the three orbits. Its
integrality is checked on the compatible outputs of the register; it is not
assumed for an arbitrary element of \(\mathbb Z^{25}\).

The three signed blocks are
\begin{equation}
 U^{(r)}=(s_r,d_{r,1}-d_{r,3},d_{r,0}+d_{r,2},d_{r,0}-d_{r,2}).
 \label{x_reciprocidad:eq:k-bloques-firmados}
\end{equation}
Integer substitution gives
\[
 (B_1,B_2,B_3)=(972,-1073,101),\qquad
 (s_1,s_2,s_3)=(2378,1406,2479),
\]
and, therefore,
\begin{align*}
 U^{(1)}&=(2378,-452,-668,-322),\\
 U^{(2)}&=(1406,998,-204,-28),\\
 U^{(3)}&=(2479,-551,-371,-997).
\end{align*}
Interleaving by columns gives
\begin{equation}
 U=(2378,1406,2479,-452,998,-551,
       -668,-204,-371,-322,-28,-997).
 \label{x_reciprocidad:eq:k-u-firmado}
\end{equation}
The presence of negative values and components greater than 999
shows that \(U\) is not a word of base-one-thousand digits. Nor is it
\(U_6\Vert U_6\), a catalog concatenation of period six. The difference
is one of domain and preserved information, not merely of values.

The constructive order thus far is
\[
 x_{\mathrm{term}}\longrightarrow\mathcal R_{12}
 \longrightarrow(b^{90},b^{120},Q_{\mathrm{TPK}})
 \longrightarrow U.
\]
The following inversion produces the dodecaphasic register. The fact that we can later recover
the channels from the dodecaphasic register constitutes a reversibility check; it does not
reverse this causal order.

\subsection{Hadamard by orbits and integer reconstruction}
\label{x_reciprocidad:subsec:k-hadamard}

In the cycle of twelve positions, step three determines three orbits:
\[
 (1,4,7,10),\qquad(2,5,8,11),\qquad(3,6,9,12).
\]
The separation of four sectors of each orbit is what, once
origin and orientation are fixed, admits the ninety-degree angular reading. The
transformation does not act on three contiguous tetrads of the natural order.

Let
\[
 H_4=\begin{pmatrix}
 1&1&1&1\\1&1&-1&-1\\1&-1&1&-1\\1&-1&-1&1
 \end{pmatrix}.
\]
If \(P_{\mathrm{orb}}\) reorders by the three indicated orbits,
we define
\[
 \mathcal H_{12}=P_{\mathrm{orb}}^{-1}
 (H_4\oplus H_4\oplus H_4)P_{\mathrm{orb}}.
\]

\begin{lemma}[Hadamard inversion and integrality]
\label{x_reciprocidad:lem:k-hadamard}
We have
\[
 H_4^2=4I_4,\qquad H_4^{-1}=\frac14H_4,\qquad\det H_4=-16.
\]
For \(u\in\mathbb Z^4\), the preimage \(H_4^{-1}u\) is integral if and
only if \(H_4u\equiv0\pmod4\).
\end{lemma}

\begin{proof}
The rows are orthogonal and have squared norm four; the matrix is
symmetric. This gives \(H_4^2=4I_4\) and the inverse. An integer elimination,
or expansion of the determinant, gives the negative sign and the value \(-16\).
The integrality condition is exactly divisibility by four
of the four coordinates of \(H_4u\).
\end{proof}

\begin{definition}[Dodecaphasic integral transformation]
The sublattice of signed registers and its integral coordinatization are
\[
 \mathcal L_H=\{u\in\ZZ^{12}:\mathcal H_{12}u\equiv0\pmod4\},
 \qquad \mathscr K(u)=\tfrac14\mathcal H_{12}u.
\]
The map \(\mathscr K:\mathcal L_H\to\ZZ^{12}\) is an isomorphism
of abelian groups, with inverse \(k\mapsto\mathcal H_{12}k\).
The canonical register \(K\) is the image of the signed register generated
by the preceding operations, with fixed phase origin and orientation.
\end{definition}

\begin{proof}
The congruence defines exactly the integrality of \(\mathscr K(u)\).
The identity \(\mathcal H_{12}^2=4I\) proves the two inverse
compositions. In particular, every \(k\in\ZZ^{12}\) has preimage
\(\mathcal H_{12}k\in\mathcal L_H\).
\end{proof}

\begin{proposition}[Canonical integral register]
\label{x_reciprocidad:prop:k-sello}
The inversion \(K^{(r)}=H_4U^{(r)}/4\) of the blocks in
\eqref{x_reciprocidad:eq:k-bloques-firmados} produces
\begin{align*}
 K^{(1)}&=(234,729,621,794),\\
 K^{(2)}&=(543,659,58,146),\\
 K^{(3)}&=(140,824,914,601).
\end{align*}
In the natural order, the dodecaphasic register is
\begin{equation}
 K=(234,543,140,729,659,824,621,058,914,794,146,601).
 \label{x_reciprocidad:eq:k-vector}
\end{equation}
It is the unique rational preimage of \(U\), and it belongs to
\(\{0,\ldots,999\}^{12}\).
\end{proposition}

\begin{proof}
The three products before division are
\begin{align*}
 H_4U^{(1)}&=(936,2916,2484,3176),\\
 H_4U^{(2)}&=(2172,2636,232,584),\\
 H_4U^{(3)}&=(560,3296,3656,2404).
\end{align*}
All coordinates are divisible by four. Dividing and undoing the
orbit permutation gives \eqref{x_reciprocidad:eq:k-vector}. Uniqueness follows from
rational invertibility, not from a search among candidates close to
a physical value.
\end{proof}

Integrality can also be formulated in terms of lattices. The
determinantal divisors of \(H_4\) are \(1,2,4,16\): the entries
have greatest common divisor one; the minors of order two are even and
one has absolute value two; those of order three are multiples of four and
one has absolute value four. It follows that
\[
 \operatorname{SNF}(H_4)=\operatorname{diag}(1,2,2,4),
\]
and for the global transformation,
\[
 \operatorname{coker}\mathcal H_{12}
 \cong(\mathbb Z/2\mathbb Z)^6\oplus(\mathbb Z/4\mathbb Z)^3.
\]
The index of the image is \(16^3=4096\). Thus not every signed integer
vector admits an integer preimage; the vector obtained satisfies the
precise congruences required by inversion. The factor \(1/4\) is
forced by the matrix and does not have the status of a free coefficient.

\subsection{Cyclic differences and uniform component}
\label{x_reciprocidad:subsec:k-transversal}

With cyclic indices modulo twelve, let
\[
 (D_3z)_m=z_m-z_{m+3},\qquad(D_4z)_m=z_m-z_{m+4},
 \qquad Q(z)=\sum_{m=1}^{12}z_m.
\]
The first operator compares sectors separated by three positions; the
second, by four. The channel names 90 and 120 here refer to
these separations in the oriented cycle, not to an identification with a
physical operator from another domain.

\begin{proposition}[Completeness of the transverse system]
\label{x_reciprocidad:prop:k-transversal}
The operator
\[
 \mathcal D=(D_3,D_4,Q):\mathbb Z^{12}\longrightarrow\mathbb Z^{25}
\]
has rank twelve and nonzero invariant factors
\[
 \operatorname{diag}(1,\ldots,1,12),
\]
with eleven units. In particular, its outputs determine a unique dodecaphasic register.
\end{proposition}

\begin{proof}
If \(D_3z=D_4z=0\), the vector is constant along both steps.
Since \(3\) and \(4\) generate the cycle modulo twelve,
\(\ker D_3\cap\ker D_4=\langle\mathbf1\rangle\). The total charge
eliminates this last freedom because \(Q(\mathbf1)=12\).

For the integral assertion, the rows of \((D_3,D_4)\) are oriented incidences
of a connected graph. A spanning tree provides eleven
unit invariant factors. Every minor of order twelve must contain the
row \(Q\). Adding the first eleven columns to the last makes the latter
zero in the incidence rows and twelve in row \(Q\). All
the maximal minors are divisible by twelve and the tree produces one of
absolute value exactly twelve. The last factor is therefore twelve.
\end{proof}

The check on \eqref{x_reciprocidad:eq:k-vector} gives
\[
 D_3K=b^{90},\qquad D_4K=b^{120},\qquad Q(K)=6263.
\]
Formulas \eqref{x_reciprocidad:eq:k-sumas-armonicas} and
\eqref{x_reciprocidad:eq:k-bloques-firmados} in turn reconstruct \(U\) from those
differences. This closes a triangle of exact representations of the
same object: transverse register, signed vector and dodecaphasic register. Equality
of results does not merge their codomains into one; it proves that
the declared maps recover the relevant information.


% BEGIN OWNER /Users/ruben/Documents/New project/output/EXTREMA_MEDIA_RAZON_RECIPROCIDAD_ES_EN_20260918/ARTICULO/ES/source/sections/registro_imagen_integral.tex
% Formalización nueva integrada en la revisión III; originales preservados.
% Propietarios: U016, ecuaciones extractor-diferencias-k y lector-armonico-*;
% registro_k.tex, subsecciones k-pi-h, k-hadamard y k-transversal.
\subsection{Integral compatibility and determination of the transverse register}
\label{x_reciprocidad:img09:seccion}

The domain of this construction is the publication of the register of an
APP--TRIT--TPK history. The evaluation of the register events,
the compatibility of its channels and the inversion of its
coordinates are distinguished. The following result gives a complete integral criterion for
the second level and an explicit factorization connecting the first to it.
The proof uses the register operators and establishes their
relations for every element of the indicated integer domains.

\subsubsection{Characterization by adjacent increments}

The indices belong to $\mathbb Z/12\mathbb Z$. Let
\[
 (Sz)_i=z_{i+1},\qquad D_3=I-S^3,\qquad D_4=I-S^4,
 \qquad Q(z)=\sum_{i=0}^{11}z_i.
\]
The transverse extractor is preserved
\[
 \mathcal D:\mathbb Z^{12}\longrightarrow\mathbb Z^{25},
 \qquad k\longmapsto(D_3k,D_4k,Q(k)).
\]
For a triple of integer coordinates $(b,c,q)$, write
\begin{equation}
 w_i=c_i-b_{i+1},\qquad
 P_0=0,\quad P_i=\sum_{j=0}^{i-1}w_j\ (1\leq i\leq12),
 \qquad T(w)=\sum_{i=0}^{11}P_i.
 \label{x_reciprocidad:img09:incremento}
\end{equation}

\begin{theorem}[Integral image of the transverse extractor]
\label{x_reciprocidad:img09:imagen}
A triple $(b,c,q)\in\mathbb Z^{25}$ belongs to
$\operatorname{im}\mathcal D$ if and only if it satisfies
\begin{align}
 b_i&=w_i+w_{i+1}+w_{i+2}&& (0\leq i<12),
 \label{x_reciprocidad:img09:relaciones}\\
 \sum_{i=0}^{11}w_i&=0,
 \label{x_reciprocidad:img09:cierre}\\
 q+T(w)&\equiv0\pmod {12}.
 \label{x_reciprocidad:img09:congruencia}
\end{align}
Its unique integer preimage is
\begin{equation}
 k_i=\frac{q+T(w)}{12}-P_i,\qquad0\leq i<12.
 \label{x_reciprocidad:img09:inversa}
\end{equation}
The twelve equations of \eqref{x_reciprocidad:img09:relaciones} and the equation \eqref{x_reciprocidad:img09:cierre} are linearly independent over $\mathbb Q$.
\end{theorem}
\begin{proof}
For an output of $\mathcal D$, one has
\[
 c_i-b_{i+1}=(k_i-k_{i+4})-(k_{i+1}-k_{i+4})=k_i-k_{i+1}.
\]
The sum of three increments gives $b_i$; their cyclic sum is zero. Moreover, $P_i=k_0-k_i$, so $q=12k_0-T(w)$. All conditions and the inverse formula are obtained.

Conversely, \eqref{x_reciprocidad:img09:congruencia} makes the formula \eqref{x_reciprocidad:img09:inversa} integer-valued; \eqref{x_reciprocidad:img09:cierre} allows it to be extended cyclically. Its adjacent differences are $w_i$. Therefore, $D_3k=b$ and
\[
 (D_4k)_i=w_i+w_{i+1}+w_{i+2}+w_{i+3}
 =w_i+b_{i+1}=c_i.
\]
The sum of the inverse gives $Q(k)=q$. The same formula proves uniqueness. For independence, the invertible integer change $(b,c)\mapsto(b,w=c-Sb)$ converts the first twelve relations into $b=(I+S+S^2)w$. Each solves for a different coordinate of $b$. The condition $\sum w_i=0$ is an additional independent relation.
\end{proof}

The rational image has dimension twelve. Its integral structure adds a congruence of order twelve: this is the constructive manifestation of the last Smith factor of $\mathcal D$. The eleven increments $w_0,\ldots,w_{10}$ and the charge $q$ are sufficient coordinates, with $w_{11}=-\sum_{i=0}^{10}w_i$ and the congruence \eqref{x_reciprocidad:img09:congruencia}. The original twenty-five coordinates preserve these same twelve degrees of freedom through thirteen relations.

\subsubsection{Commutativity with the harmonic reading}

Let $\mathcal H_{12}$ be the Hadamard transform of the register, with orbits $(r,r+3,r+6,r+9)$ for $r=0,1,2$. Its block is
\[
 H_4=\begin{pmatrix}
 1&1&1&1\\1&1&-1&-1\\1&-1&1&-1\\1&-1&-1&1
 \end{pmatrix},\qquad H_4^2=4I_4.
\]
For compatible $(b,c,q)$, define
\[
 B_r=\sum_{j=0}^3c_{r+3j},\qquad
 A_0=\frac{q+2B_0+B_1}{3},\quad A_1=A_0-B_0,
 \quad A_2=A_1-B_1.
\]
The reader $\Pi_H$ of \eqref{x_reciprocidad:eq:k-bloques-firmados} has the blocks
\[
 (\Pi_H(b,c,q))^{(r)}=
 (A_r,b_{r+3}-b_{r+9},b_r+b_{r+6},b_r-b_{r+6}).
\]

\begin{proposition}[Agreement of the two reconstructions]
\label{x_reciprocidad:img09:triangulo}
On $\operatorname{im}\mathcal D$ one has
\[
 \Pi_H\mathcal D=\mathcal H_{12},\qquad
 \mathcal D^{-1}=\tfrac14\mathcal H_{12}\Pi_H.
\]
The image of $\Pi_H$ restricted to that domain is exactly
\[
 \mathcal L_H=
 \{u\in\mathbb Z^{12}:\mathcal H_{12}u\equiv0\pmod4\}.
\]
\end{proposition}
\begin{proof}
For $k=\mathcal D^{-1}(b,c,q)$, let $a_r=\sum_{j=0}^3k_{r+3j}$. The four-position shift takes each orbit to the next, and therefore $B_r=a_r-a_{r+1}$. Together with $q=a_0+a_1+a_2$, these identities give $A_r=a_r$. On an orbit, let us write $(x_0,x_1,x_2,x_3)$ for the coordinates of $k$ and $d_j=x_j-x_{j+1}$. Then
\[
 d_1-d_3=x_0+x_1-x_2-x_3,\quad
 d_0+d_2=x_0-x_1+x_2-x_3,\quad
 d_0-d_2=x_0-x_1-x_2+x_3.
\]
They are the three remaining rows of $H_4k^{(r)}$. The inversion and the integral characterization follow from $\mathcal H_{12}^2=4I$.
\end{proof}

The condition $\Pi_H(b,c,q)\in\mathcal L_H$ controls only the
harmonic output. Membership of the twenty-five coordinates in
$\operatorname{im}\mathcal D$ additionally requires
\eqref{x_reciprocidad:img09:relaciones}--\eqref{x_reciprocidad:img09:congruencia}.
For example, adding to $c$ the vector $e_0-e_3$ preserves the three orbital
sums and leaves $\Pi_H$ unchanged. Starting from $(0,0,0)$ thus
gives a triple with $\Pi_H=0$ that violates
\eqref{x_reciprocidad:img09:relaciones}. This is an exact negative control for
channel compatibility, independent of any terminal evaluation.

\subsubsection{Determination and transport on the orbits of the register}

Theorem~\ref{x_reciprocidad:img09:imagen} determines $k$ when its channels and charge have been produced. As a condition on outputs not yet individualized, the same image is a group isomorphic to $\mathbb Z^{12}$. With only a charge $q$ fixed, the solutions form an affine class of
\[
 L_0=\{v\in\mathbb Z^{12}:\sum v_i=0\}
 =\bigoplus_{i=0}^{10}\mathbb Z(e_i-e_{11}).
\]
For example, $100\mathbf1+e_0$ and $100\mathbf1+e_1$ are two distinct registers, of charge $1201$, with components in $\{0,\ldots,999\}$. Each satisfies all image conditions and all Hadamard congruences through its own channels. These witnesses compare the compatibility equations; they are not assigned the status of terminal HMT histories.

\begin{proposition}[Cyclic covariance and stabilizer of the register]
\label{x_reciprocidad:img09:covariancia}
Let $\rho$ be window translation on a register domain $\mathscr R$ and $S$ the preceding translation on $\mathbb Z^{12}$. On the orbit of $R\in\mathscr R$, a covariant reader is determined by a vector
\[
 k\in\operatorname{Fix}(S^p),
\]
where $p\mid12$ is the minimal period of $R$ under $\rho$. The charge condition $Q(k)=q$ admits integer solutions if and only if $12/p$ divides $q$; when they exist, they form an affine class of rank $p-1$. In particular, on a free orbit, covariance and charge leave eleven degrees of freedom.
\end{proposition}
\begin{proof}
The rule $F(\rho^jR)=S^jk$ is well defined exactly when
$S^pk=k$. Such vectors repeat $12/p$ times a word of
length $p$. Their charge is $\frac{12}{p}\sum_{i=0}^{p-1}k_i$.
Divisibility and rank follow from this formula. Channel
covariance follows from $D_aS=SD_a$ and $QS=Q$.
On $\mathcal L_H$, the corresponding action is
$\widehat S_H=\mathcal H_{12}S\mathcal H_{12}/4$;
by Proposition~\ref{x_reciprocidad:img09:triangulo}, both reconstructions
intertwine these actions.
\end{proof}

If the domain carries a marked origin, the action that also transports that mark must be used. The proposition concerns this declared cyclic action; it does not presuppose the periodicity of the enriched history. Naturality with respect to truncations subsequent to the cut $108$ is obtained, for any already defined reader $F_{108}$, through $F_N=F_{108}\circ\tau_{N,108}$. This naturality preserves the initial evaluation rule; by itself it does not choose one of the possible rules on the first one hundred and eight events.

\subsubsection{Event-based evaluation and composition of contributions}

The preceding characterization provides a factorization of the extraction through two effective readings of the register:
\begin{equation}
 \omega:\mathscr R\longrightarrow
 \{w\in\mathbb Z^{12}:\sum w_i=0\},\qquad
 q_{\rm reg}:\mathscr R\longrightarrow\mathbb Z,
 \quad q_{\rm reg}(R)+T(\omega(R))\equiv0\pmod {12}.
 \label{x_reciprocidad:img09:contrato-productor}
\end{equation}
These maps must be evaluated on the events, orientation, incidence and memory produced by APP--TRIT--TPK. Once their actions on these data are given, the composition
\[
 R\longmapsto(w,q)\longmapsto
 ((I+S+S^2)w,(I+S+S^2+S^3)w,q)
 \xrightarrow{\Pi_H}u\xrightarrow{\mathcal H_{12}/4}k
\]
is determined and satisfies all inverse checks. The formula $w=c-Sb$ also allows a producer that already delivers both channels to be connected: it is a check on its output, prior to any comparison with the terminal witness.

\begin{proposition}[Accumulation of contributions and prospective uniqueness]
\label{x_reciprocidad:img09:acumulacion}
Suppose that the initial state and each event $a_t$ of the register produce, through previously fixed rules, contributions $(\delta w_t,\delta q_t)$ with
\[
 \sum_i\delta w_{t,i}=0,\qquad
 \delta q_t+T(\delta w_t)\equiv0\pmod {12}.
\]
With compatible initial conditions $(w^{(0)},q^{(0)})$, the
update
\[
 w^{(t+1)}=w^{(t)}+\delta w_t,\qquad
 q^{(t+1)}=q^{(t)}+\delta q_t
\]
produces a unique register at each stage. Its increment is
\[
 \delta k_{t,i}=
 \frac{\delta q_t+T(\delta w_t)}{12}-P_i(\delta w_t).
\]
The construction respects register concatenation, and the reading of a prefix remains unchanged under subsequent extensions.
\end{proposition}
\begin{proof}
The compatibility conditions are preserved under addition. The formula \eqref{x_reciprocidad:img09:inversa} is additive in $(w,q)$ and produces the indicated increment. Induction fixes each state successively. The sum over two segments is the concatenated sum, and a prefix receives exactly the same summands before and after extending the history.
\end{proof}

The proposition provides a sufficient criterion, not a requirement of
linearity for every TPK reader. If events are transported by nontrivial
actions, their contributions can be expressed in the terminal chart
through the corresponding affine cocycle. The cocycle equation then preserves
the same independence from the partition of the trajectory.
In both cases, the event rules and the initial state constitute the
constructive data that must come from the article: choosing them to
reproduce an expected vector would replace that provenance with another task.

The result of this development is the exact determination condition \eqref{x_reciprocidad:img09:contrato-productor}, its proved inversion and its transport identities. Its scope is the transverse register; the other HMT constructions retain their own domains and proofs.

% END OWNER


\subsection{Two rational readings: \texorpdfstring{\(\kappa_{36}\)}{kappa36} and \texorpdfstring{\(\kappa_{\mathrm{per}}\)}{periodic kappa}}
\label{x_reciprocidad:subsec:k-kappas}

Once \(K\) is constructed, the base-one-thousand chart allows two distinct
operations. The first reads a single turn; the second prolongs the dodecaphasic register
periodically. Let \(B=1000\) and define its concatenation integer
\begin{equation}
 N_K=\sum_{m=1}^{12}K_mB^{12-m}
 =234543140729659824621058914794146601.
 \label{x_reciprocidad:eq:k-numerador}
\end{equation}
The leading zeros of a block, such as that of 058, are part of its three-digit
notation. They do not change its integer value, but they do allow the twelve blocks
to be concatenated without ambiguity.

\begin{definition}[Finite and periodic readings of the dodecaphasic register]
\label{x_reciprocidad:def:k-kappas}
The reading of one turn is
\[
 \kappa_{36}=\sum_{m=1}^{12}K_mB^{-m}=\frac{N_K}{B^{12}}.
\]
The periodic prolongation is
\[
 K_j^{\mathrm{per}}=K_{1+((j-1)\bmod12)},\qquad
 \kappa_{\mathrm{per}}=\sum_{j\geq1}K_j^{\mathrm{per}}B^{-j}.
\]
\end{definition}

\begin{proposition}[Exact value and period of the prolonged dodecaphasic register]
\label{x_reciprocidad:prop:k-periodo}
We have
\begin{equation}
 \kappa_{\mathrm{per}}=\frac{N_K}{B^{12}-1},\qquad
 \kappa_{\mathrm{per}}-\kappa_{36}
 =\frac{N_K}{B^{12}(B^{12}-1)}
 =B^{-12}\kappa_{\mathrm{per}}.
 \label{x_reciprocidad:eq:k-diferencia-kappas}
\end{equation}
The infinite word of the dodecaphasic register has minimum period twelve in base one thousand and
minimum period thirty-six in base ten.
\end{proposition}

\begin{proof}
The sum of the first turn is \(N_KB^{-12}\); each new turn
multiplies its contribution by \(B^{-12}\). The geometric series produces
\(N_K/(B^{12}-1)\), and subtraction gives the second identity.

The twelve components of \(K\) are distinct. A proper period of the cycle
would identify two of them; therefore the minimum block period is
twelve. The decimal expansion is canonical, since the dodecaphasic register has no tail
of 999 blocks. If its minimum decimal period were \(d\), it would divide
36. Grouping in threes would produce a block period
\(d/\gcd(d,3)\). No proper divisor of 36 gives twelve through that
operation. Hence the minimum decimal period is 36.
\end{proof}

Both readings are rational, but they are not the same rational number. The first
ends after twelve blocks; the second repeats the turn. Moreover,
\(0<\kappa_{\mathrm{per}}-\kappa_{36}<B^{-12}\). Their closeness does not
allow one to be substituted for the other in an exact identity. In particular,
the finite closure of alpha uses \(\kappa_{36}\), whereas the
sum of the prolonged section uses \(\kappa_{\mathrm{per}}\).


% BEGIN OWNER /Users/ruben/Documents/New project/output/EXTREMA_MEDIA_RAZON_RECIPROCIDAD_ES_EN_20260918/ARTICULO/ES/source/sections/k_reversibilidad.tex
\subsection{Reversibility of the rational constant and phase change}

The rational constant \(\kappa_{\mathrm{per}}\) fully determines
the register \(K\) when the base, length and reading origin
are preserved. Its numerical function and its structural function are related
through an exact inversion.

\begin{proposition}[Integral recovery and cyclic transformation]
Let \(I(K)=\sum_{j=1}^{12}K_j1000^{12-j}\). Then
\[
 I(K)=(1000^{12}-1)\kappa_{\mathrm{per}}.
\]
The integer expansion of \(I(K)\) into twelve blocks of base \(1000\)
recovers all the \(K_j\), including the leading zeros of each block.
If \(\rho(K)=(K_2,\ldots,K_{12},K_1)\), one has
\begin{equation}
 \kappa(\rho K)=1000\kappa(K)-K_1.
 \label{x_reciprocidad:eq:k-rotacion-racional}
\end{equation}
\end{proposition}
\begin{proof}
The first identity inverts the definition of the periodic evaluation.
Iterated Euclidean division by \(1000\), with twelve fixed positions,
recovers the vector. Furthermore,
\[
 I(\rho K)=1000I(K)-K_1(1000^{12}-1).
\]
Dividing by \(1000^{12}-1\) gives \eqref{x_reciprocidad:eq:k-rotacion-racional}.
\end{proof}

Recovery therefore composes \(\kappa\mapsto K\mapsto U\)
and the transformations \(D_3K,D_4K,Q\) already constructed. The rational
evaluation preserves the signed register and the cyclic differences that
depend on it. The total depth and the transport operators of
a history remain in their own coordinates: the preceding inversion
has the integral dodecaphasic register as its domain.

Write \(\kappa_j=\kappa(\rho^jK)\), with indices modulo \(12\).
The same identity gives
\[
 K_{j+1}=1000\kappa_j-\kappa_{j+1},\qquad
 \sum_{j=0}^{11}\kappa_j=\frac{\sum_{j=1}^{12}K_j}{999}
 =\frac{6263}{999}.
\]
Phase acts exactly on the constant: its value is fixed
in the canonical reading, while its rotations follow an affine law.
The orbit sum is invariant under the choice of origin.

The subsequent incidence also has an expression in these
coordinates. The cubic threshold determines
\[
 B_K=\{j+1:1000\kappa_j-\kappa_{j+1}\ge729\}
 =\{4,6,9,10\}.
\]
The tetrad is recovered from the rational orbit because the latter
first recovers the integer components. The signed support associated
with \(\alpha\) also uses its regional registers and the boundary
normalization. The relation between coordinate and incidence is thus
expressed through identifiable operations.

This is an operational definition of the function of \(K\): it provides
a reversible integral coordinatization of the oriented register, whose
periodic evaluation and incidence functions participate in
subsequent constructions. The value, the transformation and the domain
are preserved jointly.

% END OWNER


\subsection{Periodicity and cokernel of the transport operator}
\label{x_reciprocidad:subsec:k-clase-acarreo}

The periodicity we have just proved belongs to the dodecaphasic register. It does not imply
periodicity of the regional coordinates, of the closure carries
or of alpha. It is useful to specify this separation through a second
quotient, purely integral.

Let \(S\) be the cyclic shift of twelve coordinates, with fixed
orientation, and let \(b\geq2\) be an integer base. The periodic carry
operator is \(\partial_b=S-bI\). If a periodic identity is written
\[
 A=P+E-\Phi-K+\partial_bC,
\]
then the transformations
\[
 K\longmapsto K+\partial_bh,\qquad C\longmapsto C+h
\]
preserve its right-hand side. This invariance compares representatives of a
carry class; it is not authorization to modify the dodecaphasic register already
selected by the register.

\begin{proposition}[Periodic carry quotient]
\label{x_reciprocidad:prop:k-cociente-acarreo}
With the preceding orientation,
\[
 \operatorname{coker}(S-bI)\cong\mathbb Z/(b^{12}-1)\mathbb Z.
\]
\end{proposition}

\begin{proof}
The cyclic module is presented as
\(\mathbb Z[t]/(t^{12}-1)\). Imposing \(S=bI\) is equivalent to imposing
\(t=b\), and leaves the relation \(b^{12}-1=0\). Evaluating the
coefficients at \(b\), with the chosen concatenation order, represents
the resulting class.
\end{proof}

The denominator \(B^{12}-1\) of \(\kappa_{\mathrm{per}}\) and this
quotient arise from the same cyclic length, but their objects remain
distinct: one is a rational evaluation and the other a quotient
of a carry module. The finite closure of alpha will be an oriented
chain with a right boundary. Its prolongation will transport a
compatible remote boundary, not a carry identified by decree
between the first and thirteenth positions.

The next chapter is thus prepared. The dodecaphasic register has been fixed before
solving the alpha recurrence; its Hadamard reconstruction is
integral and unique; its transverse differences recover the register; and
its two rational evaluations have separate domains and uses. None
of these facts uses alpha as input, and none yet decides the
arithmetic type of the output of the complete closure.
